\section{稠密性定理}

设 A 是一个 k-代数。若左伪模 M 是一族单左伪模的直和，则称它是半单的。

引理 1. —— 设 M 是一个半单左 A-伪模。设 B 是 \(\widetilde{A}\)-模 M 的双交换子。则 M 的每个 A-子伪模都是 M 的 B-子模。

\(\widetilde{\mathrm{A}}\)-模 M 是半单的。设 N 是 M 的一个 A-子伪模。则 N 是 M 的一个 \(\widetilde{\mathrm{A}}\)-子模。由 VIII, 第 56 页的推论 2，存在 A-模 M 的一个以 N 为像的投影子 \(p\)。由于我们有关系 \(pb = bp\)（对每个 \(b \in \mathrm{B}\)），我们得到 N 是 M 的一个 B-子模。

引理 2. —— 设 M 是一个半单左 A-伪模，并设 x 是 M 的一个元素。存在一个元素 \(a \in A\)，使得 ax = x。

设 N 是左 A-伪模 \(\widetilde{Ax}/Ax\)。它满足 \(AN = \{0\}\)。把 VIII, 第 56 页的推论 3 应用于 \(\widetilde{A}\)-模 M 与 \(\widetilde{Ax}\)，得知 A-伪模 N 是半单的。按定义，N 的每个单子伪模 S 满足 \(AS \neq \{0\}\)。于是，伪模 N 是零，并且我们得到 \(x \in Ax\)。

定理 1（雅各布森稠密性定理）. —— 设 M 是一个半单左 A-伪模。设 b 是 \(\widetilde{\mathrm{A}}\)-模 M 的双交换子的一个元素。设 \(\mathrm{F} = \{x_1, \ldots, x_n\}\) 是 M 的一个有限子集。则存在一个元素 \(a \in \mathrm{A}\)，使得 \(bx_i = ax_i\)（对每个 \(i \in [1, n]\)）。

设 B 是 \(\widetilde{\mathrm{A}}\)-模 M 的双交换子。A-伪模 \(\mathrm{M}^n\) 是半单的。设 \(\boldsymbol{x} = (x_1, \ldots, x_n) \in \mathrm{M}^n\)。由引理 2 推出 \(\boldsymbol{x} \in \mathrm{A}\boldsymbol{x}\)。\(\widetilde{\mathrm{A}}\)-模 \(\mathrm{M}^n\) 的双交换子与 B-模 \(\mathrm{M}^n\) 的位似相一致（VIII, 第 79 页，命题 2）。于是由引理 1，A-子伪模 \(\mathrm{A}\boldsymbol{x}\)（\(\mathrm{M}^n\) 的 A-子伪模）是 \(\mathrm{M}^n\) 的一个 B-子模。因此我们有包含关系 \(\mathrm{B}\boldsymbol{x} \subset \mathrm{A}\boldsymbol{x}\)。所以存在一个 \(a \in \mathrm{A}\)，使得 \(b\boldsymbol{x} = a\boldsymbol{x}\)。结论得证。
% semantic-fragments: legacy-input-seam
\relax{}
